\begin{proposition}[弱味 $CP$ 守恒的实基判据]
只考虑夸克 Yukawa 与带电流扇区。在三代模型中，
存在使 $H_u,H_d$ 同时成为实矩阵的共同左手弱基，当且仅当该扇区的 $CP$ 可以守恒。
两套质量谱均非简并时，这等价于 $J_{\rm CKM}=0$；有质量简并时应使用
\eqnrefs{eq:jarlskog-determinant-dp11} 的弱基形式。
\end{proposition}
\begin{proof}
先把 $H_u$ 对角化。非简并时保持这一对角矩阵的重相位把
$B_{ij}=(H_d)_{ij}$ 变为 $\ee^{-\ii\alpha_i}B_{ij}\ee^{\ii\alpha_j}$。
若三个非对角元均非零，可选两条相位差使 $B_{12},B_{23}$ 实；
第三条能同时实化的充要条件为
\begin{equation*}
\Im(B_{12}B_{23}B_{31})=0.
\end{equation*}
若某条边为零，其余非零边组成无闭环图，沿边依次选相位就能全实化。
再用实正交矩阵对角化实对称的 $H_d$，CKM 即可取为实矩阵。
反向结论由实矩阵的四元积虚部为零立即得到。
章末的行列式计算把这个闭环条件写成 $J_{\rm CKM}\Delta_d=0$。

若 $H_u=\operatorname{diag}(a,a,b)$，允许在前两个分量中额外作 $U(2)$ 旋转。
先用它对角化 $H_d$ 的左上 $2\times2$ 厄米子块，再用两个剩余相位把
$B_{13},B_{23}$ 分别实化；$B_{12}=0$ 已经保持，故 $H_d$ 全实。
下型简并同理。三代任意同电荷简并因此足以消去这个弱味扇区的 $CP$ 相位。
这一结论不涉及独立的强 $CP$ 角 $\bar\theta$。
\end{proof}

\begin{example}[实混合、复混合与质量简并]
比较 $V=R_{12}(\pi/6)$ 与三代标准参数化
$s_{12}=1/2$、$s_{23}=1/3$、$s_{13}=1/4$、$\delta=\pi/2$。
\end{example}
\begin{proof}[解]
第一种模型含非零代际跃迁 $|V_{us}|^2=1/4$，但所有矩阵元实，
$J=0$，所以混合本身并不要求 $CP$ 破坏。
第二种模型中
\begin{equation*}
J=\frac{5\sqrt6}{384}\simeq0.0319,\qquad
\mathscr A_\triangle=\frac{5\sqrt6}{768}.
\end{equation*}
非简并谱时它给出非零弱味 $CP$ 奇量。
若保持此参数化却令 $y_u=y_c$，则 $\Delta_u=0$，
$\det[H_u,H_d]=0$；前一证明中的简并旋转可把混合矩阵重新化实。
此时原参数化中的 $J$ 不再对全部允许的质量基变换不变，因而不能单独作物理判据。
\end{proof}
